%=====================================================
\section{Equivalence Relations and Partial Orders}
%=====================================================

%-----------------------------------------------------
\begin{frame}{Equivalence Relation}

\begin{block}{Definition}

\vspace{0.1cm}

A relation $R$ on a set $A$ is called an \textbf{Equivalence Relation} if

\[
R
\text{ is }
\boxed{\text{Reflexive}}
\land
\boxed{\text{Symmetric}}
\land
\boxed{\text{Transitive}}
\]

\end{block}

\end{frame}

%-----------------------------------------------------

\begin{frame}{Application of Equivalence Relations}

Suppose an image classifier predicts

\[
Images=
\{
Dog_1,
Dog_2,
Dog_3,
Cat_1,
Cat_2,
Car_1
\}
\]

Define
\[
xRy
\iff
\text{$x$ and $y$ belong to the same class}
\]

Equivalence Classes

\[
Dogs=
\{Dog_1,Dog_2,Dog_3\}
\]
\[
Cats=
\{Cat_1,Cat_2\}
\]
\[
Cars=
\{Car_1\}
\]

\end{frame}

%-----------------------------------------------------